\section{Conclusion}
\markright{Conclusion}

This experiment introduced the fundamental electronic components, instruments, and tools used in the Electronics Sessional. Their construction, operating principles, specifications, and laboratory applications were studied through physical observation and basic measurements.

Passive components such as resistors, capacitors, and diodes were distinguished from active devices including BJTs, MOSFETs, and operational amplifiers. The breadboard and veroboard were used for circuit prototyping, while the multimeter and oscilloscope enabled circuit testing and analysis. The soldering iron was introduced for permanent circuit assembly.

Overall, the experiment established the basic knowledge and practical skills required for subsequent laboratory experiments and future electronics engineering practice.
